\chapter{ЭКСПЕРИМЕНТАЛЬНАЯ ВАЛИДАЦИЯ}\label{ch:ch8}

В предыдущих главах были построены математическая модель системы автономных мобильных роботов с поддержкой периферийных вычислений, VCG-подобный аукционный слой распределения задач, обучаемый многоагентный контур выбора операционных режимов периферийных узлов и гибридный метод AMPLE-AMR. Также была введена масштабируемая версия C-AMPLE-AMR, основанная на кластеризации графа периферийной инфраструктуры и локализации распределения задач внутри кластеров.

Цель настоящей главы состоит в экспериментальной проверке эффективности, устойчивости и масштабируемости предложенного метода в домене роботизированного склада. Параметры складской среды и сценариев заданы в главе~\ref{ch:ch6}, а программный прототип описан в главе~\ref{ch:ch7}. Все численные результаты настоящей главы получены автоматически из воспроизводимого экспериментального контура. Таблицы и рисунки не заполнялись вручную: они сформированы на основе файлов результатов, сохранённых при полном запуске сценариев с несколькими начальными значениями генератора случайных чисел.

\providecommand{\resultplaceholder}[1]{%
\begin{center}
\fbox{\parbox{0.92\textwidth}{\textbf{Место для результата.} #1}}
\end{center}
}

\providecommand{\resulttable}[2]{%
\begin{table}[htbp]
    \centering
    \caption{#2}
    \begingroup
    \scriptsize
    \setlength{\tabcolsep}{3pt}
    \IfFileExists{#1}{\resizebox{\textwidth}{!}{\input{#1}}}{\resultplaceholder{Таблица <<#2>> будет сформирована автоматически после запуска экспериментов и сохранена в файл \texttt{\detokenize{#1}}.}}%
    \par
    \endgroup
\end{table}
}

\providecommand{\resultfigure}[3]{%
\begin{figure}[h]
    \centering
    \IfFileExists{#1.pdf}{\includegraphics[width=0.92\textwidth]{#1.pdf}}{%
    \IfFileExists{#1.png}{\includegraphics[width=0.92\textwidth]{#1.png}}{%
    \fbox{\parbox{0.88\textwidth}{\centering \textbf{Место для рисунка.}\\ \texttt{\detokenize{#3}}\\ Файл: \texttt{\detokenize{#1}}}}}}
    \caption{#2}    
    \label{#3}
\end{figure}
}

\section{Цели проверки и экспериментальные гипотезы}

Экспериментальная проверка направлена не только на сопоставление средних численных показателей, но и на проверку отдельных элементов предложенной архитектуры. Необходимо отдельно оценить вклад VCG-подобного распределения, вклад обучаемого выбора режимов периферийных узлов, вклад их совместного использования, диагностическую роль внутренних цен и целесообразность кластерного распределения.

Проверяются следующие гипотезы.

\textbf{H1. Польза VCG-подобного распределения.} При одинаковой политике выбора режимов VCG-подобный слой должен обеспечивать более высокий чистый вклад назначений в общественное благосостояние, чем простая жадная эвристика назначения.

\textbf{H2. Польза обучаемого выбора режимов.} Методы с обучаемым выбором операционных режимов периферийных узлов должны лучше адаптироваться к изменению нагрузки, очередей и сетевой обстановки, чем методы с фиксированным режимом работы узлов.

\textbf{H3. Польза гибридизации.} Полный метод AMPLE-AMR должен выигрывать не только у фиксированных базовых методов, но и у вариантов, использующих только один из двух компонентов: либо обучаемые режимы без VCG-подобного распределения, либо VCG-подобное распределение без обучаемого выбора режимов.

\textbf{H4. Диагностическая ценность внутренних цен.} VCG-платежи и оценки внешнего эффекта должны возрастать в состояниях дефицита ресурса, перегрузки узлов и конкуренции задач за один и тот же периферийный узел. Это проверяет интерпретацию платежей как внутренних теневых цен, а не как денежных переводов между независимыми владельцами.

\textbf{H5. Устойчивость к пиковым нагрузкам.} При кратковременном увеличении интенсивности поступления задач AMPLE-AMR должен снижать долю отклонённых задач и 95-й квантиль задержки по сравнению с базовыми методами.

\textbf{H6. Устойчивость к деградации связи и отказам узлов.} При ухудшении сетевой обстановки или отказе части периферийных узлов AMPLE-AMR должен сохранять более высокое общественное благосостояние и меньший хвост задержки по сравнению с методами без адаптивного выбора режимов.

\textbf{H7. Масштабируемость кластерной версии.} C-AMPLE-AMR должен снижать накладные расходы распределения в крупных конфигурациях по сравнению с глобальным AMPLE-AMR. При этом потеря общественного благосостояния относительно глобального распределения должна оставаться ограниченной.

\section{Сравниваемые методы}

В экспериментальную матрицу входят пять основных методов:
\begin{enumerate}
    \item \textbf{Fixed-Heuristic} --- фиксированный нормальный режим всех периферийных узлов и эвристическое распределение задач;
    \item \textbf{Fixed-VCG} --- фиксированный нормальный режим всех периферийных узлов и VCG-подобное распределение;
    \item \textbf{QMIX-Heuristic} --- обучаемый выбор режимов периферийных узлов и эвристическое распределение задач;
    \item \textbf{AMPLE-AMR} --- обучаемый выбор режимов периферийных узлов и VCG-подобное распределение;
    \item \textbf{C-AMPLE-AMR} --- обучаемый выбор режимов периферийных узлов и кластерное VCG-подобное распределение.
\end{enumerate}

Принципиально важно, что в наборе сравнения отсутствует вариант, в котором обучаемая политика напрямую назначает задачи на периферийные узлы без отдельного слоя распределения. Во всех экспериментах решение о режиме узла и решение о назначении задач разделены архитектурно. Это соответствует инвариантам AMPLE-AMR: обучаемая политика не изменяет заявки роботов, функции полезности, функции стоимости и правило аукционного распределения.

Для дополнительной проверки используется вариант \textbf{Random-Modes-Heuristic}, в котором режимы периферийных узлов выбираются случайно, а задачи распределяются эвристически. Этот вариант не является основным конкурентом, но позволяет проверить, связан ли выигрыш с обучением политики, а не только с самим фактом переключения режимов.

\section{Сценарии и метрики}

Для проведения экспериментальной валидации используется набор сценариев:
\begin{itemize}
    \item \texttt{stable\_warehouse\_load} --- стабильная складская нагрузка;
    \item \texttt{peak\_warehouse\_load} --- пиковая складская нагрузка;
    \item \texttt{heterogeneous\_edge\_nodes} --- разнородная периферийная инфраструктура;
    \item \texttt{network\_degradation} --- деградация связи;
    \item \texttt{edge\_node\_failures} --- отказы периферийных узлов;
    \item \texttt{scalability\_sweep} --- масштабирование числа роботов, задач и узлов;
    \item \texttt{clustered\_vs\_global} --- сопоставление глобального и кластерного распределения;
    \item \texttt{sensitivity\_operation\_modes} --- чувствительность к коэффициентам операционных режимов.
\end{itemize}

Основной интегральной метрикой является общественное благосостояние:
\[
SW(t)=
\sum_{\tau}
\sum_j
x_{\tau j}(t)
\left[
U_{i(\tau)}(\tau,n_j,t)-C_j(\tau,r_{i(\tau)},t)
\right].
\]

Для оценки качества обслуживания используются доля обслуженных задач, доля отклонённых задач, доля нарушений сроков обработки и 95-й квантиль задержки:
\[
CompletionRate = \frac{N_{completed}}{N_{generated}},
\qquad
DropRate = \frac{N_{dropped}}{N_{generated}},
\]
\[
DeadlineViolationRate = \frac{N_{violated}}{N_{completed}},
\qquad
P95Latency = \mathrm{quantile}_{0.95}\left(\{T_{\tau}:\tau\in Completed\}\right).
\]

Для оценки состояния периферийной инфраструктуры используются средняя загрузка узлов и дисбаланс загрузки:
\[
NodeUtilization = \frac{1}{k}\sum_{j=1}^{k}u_j,
\qquad
LoadImbalance = \frac{1}{k}\sum_{j=1}^{k}(u_j-\bar{u})^2.
\]

Накладные расходы принятия решения определяются как сумма времени аукционного распределения, времени вывода действия обучаемой политики и коммуникационных затрат:
\[
Overhead = T^{auction}+T^{comm}+T^{policy},
\qquad
SlotShare = \frac{T^{auction}+T^{comm}+T^{policy}}{\Delta^{slot}}.
\]

Для C-AMPLE-AMR дополнительно анализируются число кластеров, средний размер кластера, время локальных распределений, величина межкластерного среза графа и потеря общественного благосостояния относительно глобального AMPLE-AMR.

\section{Протокол экспериментов}

Полный экспериментальный прогон выполнялся по следующим правилам:
\begin{enumerate}
    \item использовались фиксированные начальные значения генератора случайных чисел $\{0,1,2,3,4\}$;
    \item для каждого масштаба складской среды политика QMIX обучалась в течение 24 эпизодов;
    \item после обучения проводилась оценка без добавления исследовательского шума;
    \item длительность управляющего слота фиксировалась равной 100~мс;
    \item длина эпизода зависела от сценария: 80 шагов для стабильной нагрузки, 90 шагов для пиковой нагрузки, 80 шагов для сценариев деградации и отказов, 60 шагов для масштабирования, 70 шагов для сравнения глобального и кластерного распределения и 70 шагов для анализа чувствительности.
\end{enumerate}

Итоговые значения в таблицах представляют собой средние по начальным значениям генератора случайных чисел. Показатели разброса дополнительно сохраняются в файлах результатов эксперимента. Число обучающих эпизодов следует рассматривать как фиксированный исследовательский бюджет обучения, достаточный для сопоставления вариантов прототипа при одинаковых условиях. Эти эксперименты не доказывают достижение глобально оптимальной политики QMIX; они проверяют воспроизводимое поведение методов при одинаковой конфигурации среды и одинаковом бюджете обучения.

\section{Поведение обучаемой политики режимов}

Перед сравнением итоговых метрик необходимо убедиться, что контур QMIX не вырождается в тривиальный фиксированный выбор одного режима. На рис.~\ref{fig:ch8_operation_mode_distribution} показано распределение операционных режимов периферийных узлов по сценариям, а на рис.~\ref{fig:ch8_learning_curves} --- динамика обучения политики.

\resultfigure{Dissertation/images/results/ch8_operation_mode_distribution}{Распределение операционных режимов периферийных узлов по сценариям}{fig:ch8_operation_mode_distribution}

\resultfigure{Dissertation/images/results/ch8_learning_curves}{Кривые обучения политики выбора операционных режимов}{fig:ch8_learning_curves}

Полученные результаты показывают, что обученная политика использует несколько режимов и меняет их распределение в зависимости от сценария. В пиковом сценарии политика почти полностью концентрируется на режимах \texttt{critical} ($74.4\%$) и \texttt{reserve} ($23.2\%$), что соответствует ожидаемому переходу к более активному использованию экспонируемой ёмкости периферийных узлов. Таким образом, обучаемый контур не сводится к случайному переключению или постоянному выбору одного режима.

\section{Стабильная складская нагрузка}

Сценарий \texttt{stable\_warehouse\_load} моделирует нормальную работу склада без отказов узлов и искусственной деградации связи. В нём проверяются гипотезы H1, H2 и H3.

\resulttable{Dissertation/tables/ch8_stable_load_results.tex}{Результаты при стабильной нагрузке}

\resultfigure{Dissertation/images/results/ch8_stable_welfare}{Общественное благосостояние при стабильной нагрузке}{fig:ch8_stable_welfare}

\resultfigure{Dissertation/images/results/ch8_stable_quality}{Доля обслуженных, отклонённых и просроченных задач при стабильной нагрузке}{fig:ch8_stable_quality}

В стабильном сценарии AMPLE-AMR превосходит Fixed-Heuristic по общественному благосостоянию на $124.0$ условных единицы. При этом Fixed-VCG практически не отделяется от Fixed-Heuristic, а глобальный AMPLE-AMR фактически совпадает с QMIX-Heuristic по основным метрикам качества. Это означает, что в данной конфигурации основной вклад в улучшение качества обслуживания даёт не VCG-подобное распределение в изоляции, а обучаемый выбор операционных режимов периферийных узлов.

\section{Пиковая нагрузка}

Сценарий \texttt{peak\_warehouse\_load} содержит кратковременный всплеск входного потока задач. В этом режиме особенно важны доля отклонённых задач, 95-й квантиль задержки и поведение операционных режимов периферийных узлов.

\resulttable{Dissertation/tables/ch8_peak_load_results.tex}{Результаты сценария пиковой складской нагрузки}

\resultfigure{Dissertation/images/results/ch8_peak_latency}{95-й квантиль задержки при пиковой нагрузке}{fig:ch8_peak_latency}

\resultfigure{Dissertation/images/results/ch8_peak_rates}{Доля обслуженных, отклонённых и просроченных задач при пиковой нагрузке}{fig:ch8_peak_rates}

В пиковом сценарии различие между методами становится существенно более заметным. Fixed-Heuristic и Fixed-VCG совпадают по основным метрикам: $SW=1608.6$, $DropRate=0.325$, $P95Latency=75.2$~мс. Напротив, QMIX-Heuristic и AMPLE-AMR увеличивают общественное благосостояние до $1819.1$, снижают долю отклонённых задач до $0.310$ и уменьшают 95-й квантиль задержки до $50.9$~мс. В абсолютных величинах выигрыш AMPLE-AMR относительно фиксированного базиса составляет $+210.5$ по благосостоянию и $-24.2$~мс по задержке.

Следовательно, H5 подтверждается: адаптивный выбор режимов повышает устойчивость системы к всплескам нагрузки. Вместе с тем AMPLE-AMR не даёт дополнительного выигрыша относительно QMIX-Heuristic в данном сценарии, поэтому вклад VCG-подобного распределения в изоляции не подтверждается.

\section{Разнородная периферийная инфраструктура}

Сценарий \texttt{heterogeneous\_edge\_nodes} проверяет способность метода учитывать различие периферийных узлов по вычислительным ресурсам, стоимости обслуживания и текущей загрузке.

\resulttable{Dissertation/tables/ch8_heterogeneous_results.tex}{Результаты сценария разнородной периферийной инфраструктуры}

\resultfigure{Dissertation/images/results/ch8_node_utilization}{Загрузка периферийных узлов в сценарии разнородной инфраструктуры}{fig:ch8_node_utilization}

В сценарии разнородной инфраструктуры AMPLE-AMR и QMIX-Heuristic снова показывают одинаковые интегральные результаты и заметно превосходят фиксированные методы: общественное благосостояние возрастает с $1207.9$ до $1303.5$, доля обслуженных задач --- с $0.760$ до $0.797$, а 95-й квантиль задержки снижается с $67.5$ до $58.9$~мс. Существенно важнее то, что обучаемые методы уменьшают среднюю загрузку узлов с $0.0225$ до $0.0184$ и снижают дисбаланс загрузки с $0.0595$ до $0.0490$.

Это означает, что обучаемый контур действительно избегает чрезмерной концентрации задач на отдельных узлах и перераспределяет нагрузку более равномерно. C-AMPLE-AMR остаётся близким по общественному благосостоянию ($1302.2$), однако уступает глобальному варианту по задержке и увеличивает накладные расходы распределения. Следовательно, в данном сценарии основное преимущество снова связано с адаптивным выбором режимов, а не с кластеризацией или самим фактом использования VCG-подобного слоя.

\section{Деградация связи и отказы узлов}

Сценарии деградации связи и отказов узлов проверяют устойчивость метода к нарушениям условий исполнения. В первом случае увеличиваются задержки или снижается качество соединения между роботами и периферийными узлами. Во втором случае часть узлов временно становится недоступной или теряет часть ресурсов.

\resulttable{Dissertation/tables/ch8_network_degradation_results.tex}{Результаты сценария деградации связи}

\resulttable{Dissertation/tables/ch8_failures_results.tex}{езультаты сценария отказов периферийных узлов}

\resultfigure{Dissertation/images/results/ch8_degradation_welfare}{Изменение общественного благосостояния при деградации связи}{fig:ch8_degradation_welfare}

\resultfigure{Dissertation/images/results/ch8_failures_quality}{Качество обслуживания при отказах периферийных узлов}{fig:ch8_failures_quality}

В сценарии деградации связи AMPLE-AMR и QMIX-Heuristic повышают общественное благосостояние с $824.2$ до $919.1$ и уменьшают 95-й квантиль задержки с $63.2$ до $58.9$~мс. При этом доля обслуженных задач почти не меняется: $0.632$ у фиксированного базиса против $0.631$ у обучаемых методов. Следовательно, выигрыш проявляется прежде всего в качестве назначений и хвосте задержки, а не в резком росте числа завершённых задач.

При отказах узлов эффект более выражен: AMPLE-AMR увеличивает общественное благосостояние с $980.2$ до $1090.9$, повышает долю обслуженных задач с $0.680$ до $0.684$ и снижает 95-й квантиль задержки с $60.2$ до $50.2$~мс. Относительно Fixed-VCG картина аналогична, поскольку Fixed-VCG и Fixed-Heuristic практически совпадают. Следовательно, H6 подтверждается: методы с адаптивным выбором режимов устойчивее методов без такого контура.

\section{Проверка интерпретации внутренних цен и внешнего эффекта}

Отдельная проверка посвящена аукционным сигналам: VCG-платежам, оценкам внешнего эффекта и показателям дефицита ресурсов. Цель проверки состоит не в том, чтобы рассматривать платежи как реальные денежные переводы, а в том, чтобы установить, являются ли они полезными внутренними диагностическими признаками.

\resulttable{Dissertation/tables/ch8_externality_diagnostics.tex}{Диагностика внутренних цен и внешнего эффекта}

\resultfigure{Dissertation/images/results/ch8_externality_vs_utilization}{Связь внешнего эффекта с загрузкой периферийных узлов}{fig:ch8_externality_vs_utilization}

\resultfigure{Dissertation/images/results/ch8_payments_by_task_class}{Внутренние цены по классам вычислительных задач}{fig:ch8_payments_by_task_class}

Данные таблицы диагностики внешнего эффекта и рис.~\ref{fig:ch8_externality_vs_utilization}--\ref{fig:ch8_payments_by_task_class} показывают, что внутренние цены реагируют на стрессовые состояния, но их диагностическая сила в текущей конфигурации остаётся умеренной. Глобальная корреляция между внешним эффектом на уровне задачи и загрузкой узла положительна, но мала: $0.021$. Средняя оценка внешнего эффекта для аукционных методов возрастает от практически нулевого уровня в стабильном сценарии ($\approx 3.8\cdot 10^{-16}$) до $2.4\cdot 10^{-4}$ при пиковой нагрузке и $9.6\cdot 10^{-4}$ в сценарии чувствительности.

Следовательно, H4 подтверждается лишь частично: сигнал существует и усиливается в условиях напряжённого обслуживания, но не является достаточно сильным, чтобы сам по себе служить надёжным индикатором перегрузки. По классам задач картина также неоднородна: универсального разделения между критическими и некритическими задачами по внутренним ценам не наблюдается. Практический вывод состоит в том, что внутренние цены полезны как вспомогательная диагностическая характеристика дефицита ресурса, однако для прямого управления обучаемым контуром их следует использовать только совместно с более устойчивыми признаками, такими как загрузка узлов и задержка.

\section{Масштабируемость и кластерная версия метода}

Сценарий масштабирования проверяет, как меняются накладные расходы распределения при увеличении числа роботов, задач, точек доступа и периферийных узлов. В этом сценарии особенно важно сравнение AMPLE-AMR и C-AMPLE-AMR. Глобальный AMPLE-AMR рассматривается как более точный, но потенциально более дорогой вариант, поскольку задача выбора распределения решается для всей инфраструктуры. C-AMPLE-AMR рассматривается как способ уменьшить вычислительные и коммуникационные накладные расходы за счёт локализации распределения внутри кластеров.

\resulttable{Dissertation/tables/ch8_scalability_results.tex}{Результаты масштабирования}

\resulttable{Dissertation/tables/ch8_clustered_vs_global_results.tex}{Сравнение глобального и кластерного распределения}

\resultfigure{Dissertation/images/results/ch8_overhead_scalability}{Накладные расходы распределения при масштабировании}{fig:ch8_overhead_scalability}

\resultfigure{Dissertation/images/results/ch8_clustered_tradeoff}{Компромисс между накладными расходами и общественным благосостоянием для C-AMPLE-AMR}{fig:ch8_clustered_tradeoff}

Результаты масштабирования не подтверждают гипотезу H7 в её сильной формулировке. Для масштабов \texttt{Warehouse-S}, \texttt{Warehouse-M}, \texttt{Warehouse-M+} и \texttt{Warehouse-L} кластерный вариант не уменьшает, а увеличивает накладные расходы относительно глобального AMPLE-AMR. Например, на \texttt{Warehouse-L} в сценарии \texttt{scalability\_sweep} суммарные накладные расходы возрастают с $2.09$ до $5.31$~мс, а в отдельном сценарии \texttt{clustered\_vs\_global} --- с $2.84$ до $7.10$~мс при одновременной потере благосостояния $5.6$.

Единственное исключение возникает на масштабе \texttt{Warehouse-XL}, где C-AMPLE-AMR уменьшает суммарные накладные расходы с $98.7$ до $50.5$~мс, то есть примерно на $48.8\%$, однако ценой потери благосостояния $32.8$. Поэтому порог, начиная с которого кластеризация начинает окупаться именно по накладным расходам, в текущем экспериментальном контуре проявляется только на самом экстремальном масштабе \texttt{Warehouse-XL} и не переносится на более умеренные конфигурации. Практически это означает, что реализация C-AMPLE-AMR требует дальнейшей доработки: либо более эффективного локального аукциона, либо иной стратегии кластеризации, прежде чем её можно будет считать общим средством масштабирования.

\section{Чувствительность к параметрам операционных режимов}

Проверка чувствительности необходима для оценки того, насколько устойчивы выводы к выбранным коэффициентам операционных режимов периферийных узлов. В данном эксперименте изменяются параметры экспонируемой ёмкости и резервирования ресурсов для режимов, используемых обучаемой политикой.

\resulttable{Dissertation/tables/ch8_sensitivity_operation_modes.tex}{Чувствительность к параметрам операционных режимов}

\resultfigure{Dissertation/images/results/ch8_sensitivity_modes}{Изменение метрик при варьировании параметров операционных режимов}{fig:ch8_sensitivity_modes}

Преимущество методов с обучаемым выбором режимов сохраняется во всех трёх профилях чувствительности. Для AMPLE-AMR общественное благосостояние составляет $915.7$ в консервативном профиле, $931.0$ в профиле по умолчанию и $951.5$ в профиле \texttt{aggressive\_bias}. Это означает, что умеренное изменение коэффициентов экспонируемой ёмкости не разрушает основной вывод о пользе адаптивного управления режимами. Более того, во всех профилях обучаемые методы превосходят как фиксированные базовые методы, так и случайное переключение режимов.

Наибольшее влияние на итоговые метрики оказывают коэффициенты, усиливающие агрессивные режимы. Профиль \texttt{aggressive\_bias} обеспечивает наибольшее благосостояние и лучшую долю обслуженных задач, тогда как профиль по умолчанию даёт минимальный 95-й квантиль задержки ($47.6$~мс). Консервативный профиль оказывается наименее выгодным по благосостоянию, хотя и не приводит к отказу метода от адаптивного поведения.

\section{Сводное сопоставление методов}

После анализа отдельных сценариев необходимо привести сводную таблицу по всем методам и сценариям. В ней отражаются общественное благосостояние, доля обслуженных задач, доля отклонённых задач, доля нарушений сроков, 95-й квантиль задержки, средняя загрузка узлов и накладные расходы распределения.

\resulttable{Dissertation/tables/ch8_overall_summary.tex}{Сводное сопоставление методов по всем сценариям}

\resultfigure{Dissertation/images/results/ch8_welfare_by_scenario}{Общественное благосостояние по сценариям}{fig:ch8_welfare_by_scenario}

\resultfigure{Dissertation/images/results/ch8_quality_by_scenario}{Качество обслуживания задач по сценариям}{fig:ch8_quality_by_scenario}

Сводное сопоставление подтверждает, что наиболее устойчивый вклад в качество обслуживания вносит именно обучаемый выбор режимов узлов. Во всех основных сценариях AMPLE-AMR превосходит Fixed-Heuristic: выигрыш по общественному благосостоянию составляет $+124.0$ при стабильной нагрузке, $+210.5$ при пиковой нагрузке, $+95.6$ в разнородной инфраструктуре, $+94.9$ при деградации связи и $+110.7$ при отказах узлов. Вместе с тем в глобальных сценариях метрики AMPLE-AMR и QMIX-Heuristic фактически совпадают, что указывает на доминирующий вклад режима управления, а не аукционного слоя распределения.

Случайное переключение режимов в отдельных случаях превосходит фиксированный режим \texttt{normal}, однако систематически уступает обучаемой политике. Следовательно, наблюдаемый выигрыш нельзя объяснить одним лишь фактом вариативности режимов. По итогам полного прогона гипотезы H2, H5 и H6 подтверждаются; H3 и H4 подтверждаются частично; H1 и H7 не подтверждаются в сильной формулировке.

\section{Выводы по главе}

В настоящей главе проведена экспериментальная проверка предложенного метода в складской AMR-среде. Эксперименты охватывают стабильную нагрузку, пиковые режимы, разнородность периферийной инфраструктуры, деградацию связи, отказы узлов, масштабирование и чувствительность к параметрам операционных режимов. Такая структура экспериментов позволила отдельно проверить вклад аукционного распределения, вклад обучаемого управления режимами и вклад кластеризации.

Полный воспроизводимый эксперимент показывает, что ключевой практический вклад метода AMPLE-AMR в текущей реализации состоит в адаптивном выборе операционных режимов периферийных узлов. Именно этот компонент обеспечивает устойчивый выигрыш по общественному благосостоянию, задержке и устойчивости к пиковым нагрузкам, деградации связи и отказам инфраструктуры.

Польза VCG-подобного слоя в изоляции от обучаемого контура в текущем наборе сценариев не подтверждается: Fixed-VCG практически не превосходит Fixed-Heuristic. Диагностическая роль внутренних цен подтверждается лишь частично: они возрастают в стрессовых режимах и могут использоваться как вспомогательные признаки дефицита ресурса, но их связь с загрузкой узлов остаётся слабой.

Переход к C-AMPLE-AMR оказывается оправдан только на экстремальном масштабе \texttt{Warehouse-XL}, где достигается сокращение накладных расходов почти вдвое. На более умеренных конфигурациях кластеризация пока не окупается. Тем самым основными ограничениями текущей реализации являются слабое разделение между эвристическим и VCG-подобным распределением при фиксированных режимах, отсутствие дополнительного выигрыша AMPLE-AMR относительно QMIX-Heuristic в глобальных сценариях и недостаточная зрелость кластерного варианта как общего механизма масштабирования.
